\documentclass[12pt]{article}
\usepackage[utf8]{inputenc}
\usepackage[T1]{fontenc}
\usepackage[margin=1in]{geometry}
\usepackage{setspace}
\usepackage{xcolor}
\usepackage[colorlinks=true, urlcolor=blue!60!black, linkcolor=blue!60!black]{hyperref}
\setlength{\parindent}{0em}
\setlength{\parskip}{0.45em}
\renewcommand{\baselinestretch}{1.0}
\begin{document}
\begin{center}
{\Large \textit{United States v. Hemani}, No. 24-1234 (2026)\par}
\vspace{0.3em}
{\normalsize Case Brief --- One-Page Edition\par}
\end{center}

\textbf{Memory Jogger.} Prosecuting a regular marijuana user for keeping a gun at home, on his drug use alone, violates the Second Amendment (\href{https://www.supremecourt.gov/opinions/25pdf/24-1234_g2bh.pdf}{Slip Op.}).

\textbf{Facts.} Ali Hemani admitted using marijuana ``about every other day'' and kept a pistol at home. He was charged only under 18 U.S.C. \S~922(g)(3), which disarms any ``unlawful user'' of a controlled substance, with no claim of addiction, intoxication, or dangerousness.

\textbf{Procedural History.} The district court dismissed the indictment; the Fifth Circuit affirmed.

\textbf{Votes.} Affirmed, 9--0 (Gorsuch for seven; Thomas and Jackson concurring; Alito and Kagan concurring in the judgment).

\textbf{Holding.} Under \textit{Bruen} and \textit{Rahimi}, the government may not disarm someone under the ``unlawful user'' provision solely for regular drug use, without pre-deprivation process or proof of incapacitation or dangerousness. The Court reserved bans on addicts and the presently intoxicated, targeted laws for drugs proven especially dangerous, \S\S~922(g)(1) and (g)(4), and prosecutions backed by individualized proof of dangerousness.

\textbf{Reasoning.} The government's ``habitual drunkard'' analogues failed on every metric: they targeted people so persistently intoxicated as to be incapable of managing their affairs, not regular drinkers; they served different purposes (labor norms, family protection, moral policing); and they operated only after judicial process. DOJ non-enforcement, state legalization, and the 2026 rescheduling of some marijuana products undercut the claim of categorical dangerousness. \emph{Separately:} Thomas doubted Congress's commerce power over \S~922(g); Jackson urged retiring \textit{Bruen} for means-end scrutiny; Alito would have affirmed only because the record showed nothing about how Hemani's use impaired him.

\textbf{Examples.} \emph{Same-Side:} a medical-marijuana user with a locked shotgun at home. \emph{Opposite-Side:} a user who is high, armed, and threatening bystanders. \emph{Fence-Sitter:} a frequent methamphetamine user with a gun but no proof of dangerousness.

\textbf{Critique.} Praised for rejecting the government's bare ``say-so''; faulted for offering little guidance and for exposing \textit{Bruen}'s incoherence.

\textbf{Key Quotation.} ``All of which leaves it awkwardly positioned to suggest that the millions of Americans who now regularly use marijuana are categorically and unusually dangerous.'' --- Justice Gorsuch, slip op. at 17.

\end{document}
